\clearpage\SourceLeaf{385}
\pdfbookmark[0]{Section VI: Various applications of the preceding investigations}{section-six}
\begin{center}\Large SECTION VI.\par\medskip\normalsize VARIOUS APPLICATIONS OF THE PRECEDING INVESTIGATIONS.\end{center}\BookRule{0.16\linewidth}

\Needspace{7\baselineskip}
\Article{308}
We have already indicated in several places how rich higher arithmetic is in truths useful in other parts of mathematics as well. We thought it worthwhile to treat separately some applications that deserve a fuller account. Our aim is less to exhaust this subject, which could easily fill several volumes, than to illustrate it with a few examples. In this section we will first discuss the decomposition of fractions into simpler ones, then the conversion of common fractions into decimals. Next we will explain a new method of exclusion for solving indeterminate equations of the second degree, and finally give new, convenient methods for distinguishing primes from composite numbers and finding the factors of the latter. In the following section we will establish the general theory of a particular kind of functions, which extends very widely throughout analysis, insofar as it is very closely connected with higher arithmetic. In particular, we will seek to advance the theory of the division of the circle, of which only the first elements have so far become known.

\Needspace{9\baselineskip}\begin{center}\textit{Decomposition of fractions into simpler ones.}\end{center}

\Needspace{7\baselineskip}
\Article{309}
\textsc{Problem.} \textit{Decompose a fraction \(\frac mn\), whose denominator \(n\) is the product of two coprime numbers \(a,b\), into two other fractions with denominators \(a,b\).}

\SourceLeaf{386}
\textit{Solution.} Let the required fractions be \(\frac xa,\frac yb\). Then we must have \(bx+ay=m\). Thus \(x\) is a root of the congruence \(bx\equiv m\; (\mathrm{mod}.\ a)\), which can be found by Section II, while \(y\)\(=\frac{m-bx}{a}\).

As we know, the congruence \(bx\equiv m\) has infinitely many roots, all congruent modulo \(a\), but just one that is positive and less than \(a\). However, \(y\) may also turn out negative. We hardly need to point out that \(y\) can also be found from the congruence \(ay\equiv m\; (\mathrm{mod}.\ b)\), and \(x\) from the equation \(x=\frac{m-ay}{b}\). — \textit{For example}, given the fraction \(\frac{58}{77}\), the number \(4\) is the value of the expression \(\frac{58}{11}\; (\mathrm{mod}.\ 7)\), so \(\frac{58}{77}\) decomposes into \(\frac47+\frac2{11}\).

\Needspace{7\baselineskip}
\Article{310}
Suppose a fraction \(\frac mn\) is given whose denominator \(n\) is the product of any number of pairwise coprime factors \(a,b,c,d\), and so on. By the preceding article it can first be decomposed into two fractions with denominators \(a\) and \(bcd\), and so on; the second can again be decomposed into two with denominators \(b\) and \(cd\), and so on; the latter again into two, and so on. Finally the given fraction takes the form \[\frac mn=\frac\alpha a+\frac\beta b+\frac\gamma c+\frac\delta d+\text{etc.}\]
Clearly we may take the numerators \(\alpha,\beta,\gamma,\delta\), and so on, positive and less than their respective denominators, except the last. Once the others have been determined, the last is no longer arbitrary, and may even be negative or greater than its denominator, since we do not assume \(m<n\). It will then usually be useful to put it in the form \(\frac\varepsilon e\mp k\), with \(\varepsilon\) positive and less than \(e\), and \(k\) an integer. Finally, \(a,b,c\), and so on, can clearly be chosen as primes or powers of primes.

\textit{Example.} The fraction \(\frac{391}{924}\), with denominator \(=4.3.7.11\), decomposes in this way into \(\frac14+\frac{40}{231}\); \(\frac{40}{231}\) into \(\frac23-\frac{38}{77}\); and \(-\frac{38}{77}\) into \(\frac17-\frac7{11}\). Thus, writing \(\frac4{11}-1\) for \(-\frac7{11}\), we get \(\frac{391}{924}=\frac14+\frac23+\frac17+\frac4{11}-1\).

\Needspace{7\baselineskip}
\Article{311}
A fraction \(\frac mn\) can be put in the form \(\frac\alpha a+\frac\beta b+\text{etc.}\mp k\) in \textit{only one way}, with \(\alpha,\beta\), and so on, positive and less than \(a,b\), and so on. For suppose \[\frac mn=\frac\alpha a+\frac\beta b+\frac\gamma c+\text{etc.}\mp k=\frac{\alpha'}a+\frac{\beta'}b+\frac{\gamma'}c+\text{etc.}\mp k'\]
\SourceLeaf{387}
with \(\alpha',\beta'\), and so on, likewise positive and less than \(a,b\), and so on. Then necessarily \(\alpha=\alpha',\ \beta=\beta',\ \gamma=\gamma'\), and so on, and \(k=k'\). Multiplying by \(n=abc\), and so on, clearly gives \(m\equiv\alpha bcd\text{ etc.}\equiv\alpha'bcd\text{ etc.}\; (\mathrm{mod}.\ a)\). Since \(bcd\), and so on, is coprime to \(a\), necessarily \(\alpha\equiv\alpha'\), and hence \(\alpha=\alpha'\). Likewise \(\beta=\beta'\), and so on, from which \(k=k'\) also follows at once. Now we are entirely free to choose which denominator's numerator to calculate first. So \textit{all} the numerators can clearly be found in the same way as \(\alpha\) in the preceding article: \(\beta\) from the congruence \(\beta acd\text{ etc.}\equiv m\; (\mathrm{mod}.\ b)\), \(\gamma\) from \(\gamma abd\text{ etc.}\equiv m\; (\mathrm{mod}.\ c)\), and so on. The sum of all the fractions found in this way either equals the given fraction \(\frac mn\), or differs from it by an integer \(=k\). This also gives us a check on the calculation. Thus, in the example of the preceding article, the values of the expressions \(\frac{391}{231}\; (\mathrm{mod}.\ 4)\), \(\frac{391}{308}\; (\mathrm{mod}.\ 3)\), \(\frac{391}{132}\; (\mathrm{mod}.\ 7)\), \(\frac{391}{84}\; (\mathrm{mod}.\ 11)\) immediately give the numerators \(1,2,1,4\) corresponding to the denominators \(4,3,7,11\), and the sum of these fractions is found to exceed the given fraction by one.
